% 국문 초록
\chapter*{국문 초록}
\phantomsection\addcontentsline{toc}{chapter}{국문 초록}

\begin{center}
{\large\bfseries \ThesisTitleKo\par}
{\normalsize --- \ThesisSubtitleKo\ ---\par}
\vspace{2mm}
{\normalsize \AuthorKo\quad \DepartmentKo\ \UniversityKo\ 대학원\par}
\end{center}
\vspace{4mm}

갑상선 결절의 초음파 판독에서 결절의 경계(margin)는 악성도를 가늠하는 독립적인 채점 항목이다.
따라서 자동 분할 모델에 요구되는 것은 영역이 대체로 맞는 것을 넘어 경계가 정확한 것이며,
검진 현장이나 휴대형 초음파처럼 GPU가 없는 환경에서도 동작해야 실제로 쓰일 수 있다.
본 연구는 YOLOv11-seg를 기반으로, 두 공개 데이터셋을 혼합한 학습, 추론 비용을 늘리지 않는
경계 정교화, 사후 양자화를 통한 경량화를 하나의 파이프라인에서 검증하였다.

먼저 TN3K와 DDTI를 합친 4,130장(결절 4,457개)에 대해 출처 계층화 5-fold 교차검증과
학습에 쓰지 않은 TN3K test 614장의 held-out 평가, 그리고 한 데이터셋으로 학습해 다른 데이터셋에서
평가하는 LODO(leave-one-dataset-out) 실험을 설계하였다. 이 조건에서 YOLOv11n-seg(nano)는
small 모델보다 파라미터가 3.5배 적음에도 Mask \mAP{} 0.6272로 small(0.6239)과 동등하여
경량 모델을 기준 모델로 선택하였다.

다음으로 분할 헤드의 프로토타입 특징에서 갈라지는 학습 전용 경계 분기(\BRM{})와,
정답 마스크의 형태학적 경계 밴드에 가중치를 주는 경계 인식 손실을 제안하였다.
경계 분기는 4,657개의 파라미터(전체의 0.16\%)만 추가하며 추론 시에는 실행되지 않으므로
배포되는 연산 그래프는 베이스라인과 동일하다(9.73 GFLOPs). 손실 계수를 스윕으로 정한 뒤
held-out 평가에서 Dice는 0.8195에서 0.8300으로, HD95는 영상 대각선 대비 9.628\%에서 8.378\%로
13.0\% 개선되었고, Boundary IoU도 0.4358에서 0.4414로 향상되었다. 영상 단위 대응표본 검정(n=614)에서
세 지표 모두 유의하였다(Wilcoxon $p$ = 0.014, 0.0045, 0.011). 반면 교차 데이터셋 조건에서는 개선이
유지되지 않았으며, 3개 시드로 반복한 결과 악화 역시 통계적으로 확인되지 않아
일반화 여부는 판정할 수 없었다. 단일 시드 결과만으로는 큰 악화처럼 보였던 차이가
시드 반복 후 실행 간 변동 범위 안으로 들어왔다는 점도 함께 보고한다.

마지막으로 제안 모델을 ONNX FP32·FP16, ONNX Runtime 정적 INT8, OpenVINO INT8로 변환하여 비교하였다.
OpenVINO INT8 모델은 3.71 MB, CPU 지연 4.10 ms로 ONNX FP32(11.63 MB, 62.82 ms) 대비 15.3배 빨랐고,
Dice 0.8290으로 원본(0.8300)과 사실상 같은 정확도를 유지하였다. 같은 INT8이라도 ONNX Runtime에서는
검출 헤드를 부동소수점으로 남겨야 했기 때문에 44.00 ms에 그쳐, 양자화 효과가 변환 도구의
연산 선별 방식에 크게 좌우됨을 확인하였다.

\vspace{6mm}
\noindent\textbf{주요어}\quad 갑상선 결절, 초음파 영상, 인스턴스 분할, YOLOv11, 경계 정교화,
경계 인식 손실, 모델 경량화, 사후 양자화, 교차 데이터셋 일반화

\vspace{2mm}
\noindent\textbf{학\quad 번}\quad \fillin{○○○○-○○○○○}
